\documentclass[a4paper,11pt]{article}

\usepackage{revy}
\usepackage[utf8]{inputenc}
\usepackage[T1]{fontenc}
\usepackage[danish]{babel}

\revyname{Matematikrevyen}
\revyyear{2026}
\version{0.1}
\eta{$2$ minutter}
\status{Færdig}

\title{Dobbelttjek}
\author{Niklas '21}

\begin{document}
\maketitle

\begin{roles}
\role{S1}[Skuespiller] "Studerende"
\role{S2}[Skuespiller] "Studerende"
\end{roles}

\begin{sketch}
\scene{S1 og S2 sidder ved et bord or arbejder. S1 har en lommeregner.}
\says{S1}[Koncentreret] ... og som det sidste... cosinus af 52 er 0.62... Så er afleveringen færdig.
\says{S2} Perfekt, lad mig bare lige...
\scene{S2 tager egen lommeregner frem. S1 stopper S2, skubber lommeregneren på bordet og kigger lidt forarget.}
\says{S1} Hvad laver du?
\says{S2}[Løfter lommeregneren igen] Jeg ville bare lige tjekke at det er rigtigt?
\says{S1} Jamen jeg har jo lige regnet det.
\scene{S2 rækker ud efter lommeregneren, men S1 stirrer intest på S2, så S2 lægger den ned igen og kigger akavet frem.}
\says{S2} ... Det kan jo ikke skade at dobbeltjekke det.
\says{S1} Men det ville heller ikke blive bedre... Jeg har jo regnet det.
\scene{Der er lang akavet stilhed.}
\says{S2} Det er bare... jeg har virkelig brug for 12 i det her fag.
\says{S1}[Irriteret] Altså tror du ikke at jeg kan bruge en lommeregner? Tror du jeg er rus?
\says{S2} Jamen... rus er de eneste der stadigvæk kan bruge lommeregner?
\says{S1} Så du kalder mig rus?? så du siger at jeg tager fejl?
\says{S2} Nej jeg tror ikke du tager fejl, det er ba-
\says{S1}[mere irriteret] Jamen det er jo det du siger. Hvorfor ville du spilde din tid? Hvorfor ville tanken overhovedet strejfe dig hvis ikke det var fordi du har tvivl i dit hjerte... omkring mine tekniske evner...
\scene{S1 er helt forpustet efter sin rant. S1 slapper mere af. S2 ser mere utilpas ud og kigger på sin lommeregner. Det opdager S1 og begynder igen.}
\says{S1} Du ville jo ikke vise kontrapositionen af en biimplikation der allerede er vist. Det er gjort... det er done...
\says{S2} Men okay hvornår er din lommeregner sidst blevet kalibreret?
\says{S1}[Intenst helt oppe i ansigtet af S2] Tror du jeg render rundt med en ukalibreret lommeregner? Hvad tager du mig for?
\scene{Lang akavet pause}
\says{S2} Men altså er du sikker?
\says{S1}[Holder sin lommeregner op] Det her er en vaskeægte japansk PI314, den har aldrig regnet noget forkert. Hvem er du, der kan stille spørgsmålstegn ved så fint et instrument?
\says{S1} Jeg har fået nok, nu afleverer jeg altså.
\scene{S1 begynder at taste på sin computer}
\says{S2} Nej, stop!
\scene{S1 prøver at uploade imens S2 prøver at stoppe dem. De slås mildt, men S1 klikker enter meget tydeligt}
\says{S1} Så er den sendt afsted!
\scene{S2 ser meget utilfreds ud. Efter noget tid kommer der en lyd fra computeren.}
\says{S1} Wow, vi har allerede fået afleveringen tilbage.
\says{S2}[Tager computeren] Giv mig den der.
\says{S2} Dumpet? Hvordan er den dumpet?? Lad mig tjekke feedback... "For mange regnefejl"...
\scene{S2 kigger surt på S1. S1 trykker hektisk på sin lommeregner. S1 slår sig selv i panden}
\says{S2} Hvad har du gjort..?
\says{S1} Min lommeregner var sat til degrees.
\scene{Hvis man kunne time det, så lys ned imens man kan se S2 række ud efter S1 for at kvæle dem.}



\end{sketch}
\end{document}